\chapter{Conclusions and Future Work}
\label{ch:conclusions_future_work}

\section{Conclusions}
\label{sec:conclusions}

The developed application meets the main objective of processing, analysing
and comparing several transient cases associated with the same
\acrshort{esatan-tms} thermal model. The initial steady-state tool has been
extended without removing its previous inspection capabilities, providing a
common environment for examining complete time series and results at selected
instants.

A representative spacecraft thermal model has been developed to support the
implementation and evaluation of these functions. The nominal, hot and cold
cases retain the same geometry, node hierarchy, materials and thermal
connections, while the environmental conditions and internal dissipations are
modified. This common basis allows equivalent groups, nodes and exported
result fields to be compared directly throughout the simulated interval.

The application can load several compatible result files, preserve their case
identity and apply the same selection to all of them. Transient temperatures,
thermal loads, conductive and radiative heat flows, interface temperature
differences and time-integrated energy transfers can be represented using
superposed curves. The \textit{Case difference} option complements these
absolute results by calculating the separation between a reference case and
one or more comparison cases. Together, both representations preserve the
physical operating levels and make smaller differences or non-coincident
orbital events easier to identify.

Additional calculations extend the imported \acrshort{esatan-tms} results.
Temperature changes can be evaluated over selected time windows, critical
interfaces are identified from explicit thermal connections and summary tables
condense the main extrema and their occurrence times. Analyses and their
configurations can also be saved, restored and exported, reducing the need to
repeat the complete selection process when the same result must be examined
again.

The thermal post-processing is complemented by an electrical analysis based on
the selected solar-panel and internal-load results. The absorbed solar load is
used to estimate incident and converted electrical power, while the selected
internal dissipations provide an approximation of the electrical load demand.
These profiles supply a separate battery model that calculates state of charge,
terminal voltage, current, accepted and curtailed power, conversion and
resistive losses, and temperature-dependent operating restrictions. All these
results are obtained during post-processing and do not modify the original
thermal solution.

The comparison of the three cases confirms that the highest absolute
temperature does not necessarily coincide with the largest temporal variation
or the greatest local gradient. The hot case produces the highest equipment
temperatures and the largest OBC--radiator interface difference, whereas the
cold case produces the lowest battery temperature and activates the charging
derating. Differences in eclipse timing also affect solar generation, load
support and the final battery state.

The hot case generates the greatest electrical energy but finishes with the
lowest state of charge because its late eclipse and higher demand produce the
strongest final discharge. The cold case receives less solar energy and
remains thermally restricted, although its lower load prevents unserved
demand. These results show that battery operation cannot be determined from
the generated or demanded energy alone, since their temporal distribution and
the active operating limits are also relevant.

The resulting tool therefore connects transient thermal inspection,
multi-case comparison and a simplified electrical interpretation within the
same workflow. It does not replace the thermal solver or constitute a complete
spacecraft electrical-power-system model. Its purpose is to organise the
exported results, reproduce consistent analyses and reveal the relations
between environmental conditions, thermal response, electrical generation and
battery operation. The modular structure also provides a suitable basis for
further extensions without requiring the existing analysis areas to be
redesigned.


\section{Future Work}
\label{sec:future_work}


The following lines extend the current physical models, analysis capabilities
and software architecture without changing the scope of the completed
implementation.

\begin{itemize}[
    leftmargin=*,
    itemsep=0.9em,
    topsep=0.5em,
    parsep=0pt
]
    \item \textbf{Battery-model extension.}
    The current implementation is a configurable post-processing approximation
    based on a generic piecewise voltage curve, separate charge and discharge
    resistances and a simplified CC--CV strategy. Greater physical fidelity
    would require voltage--SOC curves calibrated for a specific cell, different
    charge and discharge characteristics, temperature-dependent resistance and
    usable capacity, hysteresis and dynamic RC branches. Ageing,
    cycle-dependent capacity loss and validation against experimental or
    manufacturer data would also allow the calculated voltage, current and
    state of charge to be evaluated for a defined battery technology.

    \item \textbf{Bidirectional thermal--electrical coupling.}
    The thermal and electrical calculations currently follow a one-directional
    sequence: thermal results determine solar generation, load demand and
    battery operation, but the calculated electrical losses are not returned
    to the thermal model. A bidirectional procedure could convert battery
    resistive losses, conversion losses and operating-mode changes into
    time-dependent heat loads for a subsequent \acrshort{esatan-tms} solution.
    An iterative exchange between both models would provide a more consistent
    thermo-electrical response without requiring the post-processing
    application to replace the thermal solver.

    \item \textbf{Comparison between different thermal-model architectures.}
    Multi-case comparison is presently restricted to result files that share a
    common node hierarchy. A further extension could compare successive
    versions or different architectures of the same thermal model by defining
    correspondences between equivalent groups and components. The comparison
    would then identify common elements, warn about unmatched regions and
    normalise selected results by properties such as area, mass or thermal
    capacitance when a direct nodal subtraction is not meaningful.

    \item \textbf{Alignment by orbital events.}
    Temporal alignment by eclipse entry, eclipse exit or orbital phase would
    permit comparisons between simulations with different starting instants or
    output-time vectors. This option would separate the effect of the physical
    event from a displacement in the simulation timeline and facilitate the
    comparison of cases with different eclipse durations.

    \item \textbf{Sensitivity and uncertainty analysis.}
    The existing multi-case framework provides a basis for sensitivity and
    uncertainty studies. Sets of cases could be generated by varying contact
    conductances, thermo-optical properties, internal dissipations, orbital
    conditions or electrical parameters. Automated envelopes, uncertainty
    bands and rankings would identify the inputs with the greatest influence on
    each component.

    \item \textbf{Automatic critical-event detection.}
    The same process could detect temperature-limit violations, maximum case
    differences, rapid thermal variations, changes in heat-flow direction,
    minimum state of charge, active derating and periods of unserved load. A
    common event table could then identify the affected case, element, time and
    result without requiring every curve to be inspected individually.

    \item \textbf{Performance and scalability.}
    Larger thermal models and broader collections of transient cases would
    benefit from selective loading of required result fields, deferred
    processing, reuse of intermediate calculations and parallel execution
    between cases. These improvements could reduce memory use and response time
    when processing long simulations, finer temporal discretisations or a
    large number of result files.

    \item \textbf{Batch processing and automatic report generation.}
    A batch mode could apply a saved analysis configuration to an entire folder
    of cases and export the corresponding plots, tables and numerical
    summaries. Report templates for LaTeX, PDF or spreadsheet formats would
    extend this process towards reproducible and partially automated result
    documentation.

    \item \textbf{Persistent-storage improvements.}
    Metadata, widget configurations and generated files are currently
    coordinated between a local database and output folders. A unified
    transactional layer could prevent incomplete records, detect orphaned
    files and avoid unnecessary duplication. Portable analysis packages would
    allow complete configurations to be exported, transferred and restored on
    another installation. Version information and migration rules would
    preserve compatibility when the application structure or available
    controls change.

    \item \textbf{Dedicated web architecture.}
    The interface may evolve from Streamlit towards a separate frontend based
    on HTML, CSS and JavaScript, connected to a Python backend through an
    application programming interface. This architecture would offer greater
    control over layout, asynchronous calculations, user sessions and remote
    deployment while retaining the calculation modules as an independent
    backend.

    \item \textbf{Automated testing and validation.}
    Unit, integration and regression tests could check file loading, sign
    conventions, numerical integration, multi-case differences, battery
    calculations and saved-analysis restoration after each software change.
    These tests would also provide a stable basis for introducing new result
    formats, analysis modules and physical models.
\end{itemize}
